% Use llncs and set the depth for the table of contents to 2 (won't show sub-subsections)
\documentclass[runningheads, a4paper]{llncs}
\setcounter{tocdepth}{2}

% \usepackage{geometry}
% \geometry{}

% Biblography

% \usepackage[sorting=nyt,style=apa]{biblatex}
% \usepackage[style=authoryear]{biblatex}
% \let \cite \textcite 
% \usepackage{biblatex}
\usepackage[sorting=nyt,style=apa]{biblatex}
% \usepackage{harvard}
% \usepackage[style=authoryear-ibid]{biblatex})
\addbibresource{list_included.bib}
\addbibresource{list_reading.bib}
% \bibliographystyle{src/splncs04}
% \bibliography{list_included, list_reading}

% Various packages accrued over previous projects
\usepackage{graphicx}
\usepackage[]{hyperref}
\usepackage[x11names]{xcolor} % xcolor.pdf/4.3 for palettes
\usepackage[]{todonotes} % This provides text boxes for notes
\usepackage[]{subcaption} % For subfigures environments https://tex.stackexchange.com/a/282878
\usepackage[]{caption}
\usepackage[]{pgfgantt}
\usepackage[]{changepage} % This provides adjustwidth
\usepackage{cleveref}

\usepackage[]{lipsum}

% From https://tex.stackexchange.com/a/181233
\captionsetup[subfigure]{subrefformat=simple,labelformat=simple}
\renewcommand\thesubfigure{(\alph{subfigure})}
% \usepackage[title,toc,titletoc,header]{appendix} % Nice appendix sections



\hypersetup{
    colorlinks,
    allcolors=SteelBlue3, % Setup the colour of the reference links
    % Using '.' sets it to whatever the current text colour is, requires "xcolor"
    urlcolor=blue,
}

% If you use the hyperref package, please uncomment the following line
% to display URLs in blue roman font according to Springer's eBook style:
% \renewcommand\UrlFont{\color{blue}\rmfamily}
\urlstyle{rm} % The above is broken, this works (along with the blue from above)
% Although it seems that the font doesn't change or is the same


% Macros

\let \cite \parencite
\newcommand \todoask[1]{ \todo[color=MediumOrchid1]{#1} }
\newcommand \todoaskk[1]{ \todo[color=MediumPurple1]{#1} }
\newcommand \nt[1] { \textcolor{DeepPink2}{ \textless{}{#1}\textgreater } }
\newcommand \rf { \textcolor{DeepPink2}{ \textless{}REF?\textgreater } }
\newcommand \rfq { \textcolor{DeepPink2}{ \textless{}ref maybe?\textgreater } }
\newcommand \hlchanged[1] { \colorbox{LightSteelBlue1}{\parbox{12cm}{ {#1} }} }
\newcommand \hldeleted[1] { \colorbox{Firebrick1}{\parbox{12cm}{ {#1} }} }


% This macro was taken from here:
% https://tex.stackexchange.com/questions/173820/how-to-put-a-matrix-of-images#173826
\newcommand{\subf}[2]{%
  {\small\begin{tabular}[t]{@{}c@{}}
  #1\\#2
  \end{tabular}}%
}

% This macro was taken from here:
% https://tex.stackexchange.com/questions/44618/dynamically-count-and-return-number-of-words-in-a-section
%\newcommand\wordcount{
%    \immediate\write18{texcount -sub=section \jobname.tex  | grep "Section" | sed -e 's/+.*//' | sed -n \thesection p > 'count.txt'}
%(\input{count.txt}words)}


\begin{document}

    \title{\LaTeX \hspace{0.09em} Simple (\& Hacky) Template} %\thanks{Supported by organization x.}}
    \subtitle{This is a subtitle}
    \titlerunning{Shorter title for Simple (\& Hacky) \LaTeX \hspace{0.09em} Template}
    
    \author{
        Author First \inst{1,2} % \and % \orcidID{0000-1111-2222-3333}
    }
    % \and
    %     Second Author\inst{2,3}\orcidID{1111-2222-3333-4444} \and
    %     Third Author\inst{3}\orcidID{2222--3333-4444-5555}}
    % %
    \authorrunning{Author First}
    % First names are abbreviated in the running head.
    % If there are more than two authors, 'et al.' is used.
    %
    \institute{Institution; \email{email@domain.ac.uk}
    \and
    Institution 2; \email{email2@domain.ac.uk}}
    % \and
    %     Springer Heidelberg, Tiergartenstr. 17, 69121 Heidelberg, Germany
    %     \email{lncs@springer.com}\\
    %     \url{http://www.springer.com/gp/computer-science/lncs} \and
    %     ABC Institute, Rupert-Karls-University Heidelberg, Heidelberg, Germany\\
    %     \email{\{abc,lncs\}@uni-heidelberg.de}}
    %
    \maketitle              % typeset the header of the contribution
    %
    \begin{abstract}
        The abstract should briefly summarize the contents of the paper in
        15--250 words.

        \keywords{\LaTeX \and Template \and Another keyword.}
    \end{abstract}

    \tableofcontents
    \listoffigures
    \listoftables

    \newpage
    
    \section{Boilerplate}

        Here we have a bunch of lorem ipsum text, Fig. \ref{fig:antarctic} showing an antarctic scene (no bears! $\leftarrow$ this is a joke about etymology, see \textcite{antarcticEtymonline,articEtymonline}), and Figure \ref{fig:toriiAndRiver} containing \ref{fig:toriiAndRiver:torii} or \ref{sub@fig:toriiAndRiver:torii} which depicts a torii and \subref{fig:toriiAndRiver:river} depicting a river scene.

        \lipsum[1]

        \begin{figure}[h]
            \begin{center}
                \begin{scriptsize}
                    \vspace{-.5cm}
                    \begin{adjustwidth}{-0cm}{-0.cm}
                        \begin{center}
                            \includegraphics[width=0.5 \linewidth]{images/antarctic-svgrepo-com.svg.eps}
                        \end{center}
                    \end{adjustwidth}
                \end{scriptsize}
                \vspace{-.0cm}
                \caption{Antarctic Scene}{A vector graphic scene from \href[]{https://www.svgrepo.com/}{svgrepo.com} picturing an antarctic research station}
                \label{fig:antarctic}
            \end{center}
            \vspace{-0.6cm}
        \end{figure}


        \begin{figure}
            \begin{center}
                \begin{subfigure}[h]{0.4 \linewidth}
                    \includegraphics[width=\linewidth]{images/shinto-japan-svgrepo-com.svg.eps}
                    \caption{Shinto torii scene}
                    \label{fig:toriiAndRiver:torii}
                \end{subfigure}
                \hfill
                \begin{subfigure}[h]{0.4 \linewidth}
                    \includegraphics[width=\linewidth]{images/landscape-river-svgrepo-com.svg.eps}
                    \caption{River landscape scene}
                    \label{fig:toriiAndRiver:river}
                \end{subfigure}
                \caption{Torii and River Scenes}{Two side-by-side vector graphics converted to Encapsulated PostScript format from \href[]{https://www.svgrepo.com/}{svgrepo.com} showing both \subref{fig:toriiAndRiver:torii} a torii and \subref{fig:toriiAndRiver:river} a river scene}
                \label{fig:toriiAndRiver}
            \end{center}
            \vspace{-0.6cm}
        \end{figure}

        \lipsum[2]

        \lipsum[3]

    \section{Table Section}
        
        \begin{table}
            \begin{adjustwidth}{-1cm}{-1cm}
                \begin{tabular}{|| l | c | r | p{8cm} ||}
                    Left-aligned & Centred & Right-aligned & 8cm Paragraph \\
                    \hline
                    Nunc & eleifend & consequant & Donec bibendum quam in tellus. \\
                    Nam & ipsum & ligula & Quisque ullamcorper placerat ipsum. Cras nibh. Morbi vel justo vitae lacus tin-cidunt ultrices. Lorem ipsum dolor sit amet, consectetuer adipiscing elit. \\
                    Proin & tempus & nibh & Aenean placerat. \\
                \end{tabular}
            \end{adjustwidth}
                
            % Give a bit of a margin to the table
            \vspace{0.25cm}

            \begin{center}
                \caption{A Table}{This is a table, with a short caption, this caption and a label}
                \label{tab:aTable}
            \end{center}
        \end{table}

    \section{Subsection Section}

        \lipsum[4]
        
        \subsection{A Very Important Subsection}

            \lipsum[5]

            \subsubsection{An even more important sub sub-section}

                \lipsum[6]

                \vspace{-0.5cm} % You can also overlap sections if you don't like your reader
                Note that sub-subsections aren't included in the TOC with a TOC depth of "2", "3" is needed. See Table \ref{tab:aTable}.


    %\pagebreak
    \newpage
    
    %
    % ---- Bibliography ----
    %
    % BibTeX users should specify bibliography style 'splncs04'.
    % References will then be sorted and formatted in the correct style.
    %

    % \bibliographystyle{src/splncs04}
    \printbibliography

    
\end{document}
